\section{Triangle shape-space stability and the equilateral approximation}
\label{sec:triangle_stability}

\subsection{Question}

Sections~\ref{sec:efimov}--\ref{sec:wavefunction} solve the Schr\"odinger
equation along the equilateral \emph{breathing mode} $d_{12}=d_{13}=d_{23}=d$.
Is the equilateral configuration actually a stable energy minimum in shape
space, or is it a saddle point?  The answer determines whether the 1D
reduction is exact or approximate.

\subsection{Parameterisation of shape deformations}

We parameterise deviations from equilateral by two dimensionless parameters:
\begin{equation}
  d_{12} = d(1+\varepsilon_a),\quad
  d_{13} = d(1+\varepsilon_b),\quad
  d_{23} = d.
\end{equation}
The ZP term uses $\bar{d}=(d_{12}+d_{13}+d_{23})/3$.
At $(\varepsilon_a,\varepsilon_b)=(0,0)$ the triangle is equilateral.

By $S_3$ symmetry, the full Hessian in the 3-distance space is
$H_{ij}^{(3)} = \alpha\,\delta_{ij} + \beta\,(1-\delta_{ij})$,
with two eigenvalue families:
\begin{align}
  \text{Breathing mode:} &\quad \lambda_{\rm br} = \alpha + 2\beta
    \quad (\text{equilateral scaling}), \\
  \text{Shape modes (2-fold degenerate):} &\quad \lambda_{\rm sh} = \alpha - \beta.
\end{align}

\subsection{Numerical results at \texorpdfstring{$C=C_{\rm Pl}$}{C=C\_Pl},
  \texorpdfstring{$m_{\rm sp}=0.1$}{m\_sp=0.1}}

\paragraph{Classical equilibrium ($d_{\rm cl}=1.37\,l_{\rm Pl}$):}

\begin{equation}
  H = \begin{pmatrix}
    -0.0144 & +0.1026 \\
    +0.1026 & -0.0146
  \end{pmatrix}\;\text{E}_{\rm Pl}
  \quad\Rightarrow\quad
  \lambda_{\rm sh} = -0.117,\quad \lambda_{\rm br}^{(2)} = +0.088\;\text{E}_{\rm Pl}.
\end{equation}

Both eigenvalues of the reduced $2\times2$ matrix are identified: the
\emph{antisymmetric} shape mode (one bond shortens, another lengthens)
has $\lambda_{\rm sh}=-0.117<0$ — the equilateral is a \textbf{saddle point}
at its own classical minimum.

\begin{table}[h]
\centering
\caption{Shape-mode Hessian eigenvalues vs $d$ at $C_{\rm Pl}$, $m_{\rm sp}=0.1$.
Negative eigenvalue indicates local instability in that direction.}
\label{tab:shape_hessian}
\begin{tabular}{rlll}
\hline
$d$ $[l_{\rm Pl}]$ & $\lambda_1$ & $\lambda_2$ & Classification \\
\hline
$1.00$ & $-0.160$ & $+0.278$ & Saddle \\
$1.37$ (class.\ min.) & $-0.117$ & $+0.088$ & Saddle \\
$1.84$ & $0$ & $-0.030$ & Transition \\
$2.00$ & $-0.081$ & $-0.016$ & Fully unstable \\
$4.52$ (quant.\ $\langle d\rangle$) & $-0.064$ & $-0.037$ & Fully unstable \\
$8.00$ & $-0.051$ & $-0.022$ & Fully unstable \\
\hline
\end{tabular}
\end{table}

\noindent
The stability transition (where the larger shape eigenvalue crosses zero)
occurs at $d^*\approx1.84\,l_{\rm Pl}$.  For $d>d^*$, the equilateral
is a \emph{local maximum} in both shape directions — the potential energy
decreases in all shape-deformation directions.

\paragraph{At the quantum mean position $\langle d\rangle=4.52\,l_{\rm Pl}$:}
Both eigenvalues are negative ($\lambda_1=-0.064$, $\lambda_2=-0.037$),
confirming the instability extends through the entire classically accessible
region of the quantum wavefunction.

\subsection{Physical interpretation}

The shape instability has a clear physical origin: at $C=C_{\rm Pl}$,
the pair interaction $V_2\propto e^{-m_{\rm sp}d}/d$ is non-negligible.
An isosceles deformation creating one shorter bond $d_{\rm short}<d$ and
one longer bond $d_{\rm long}>d$ with $d_{\rm short}+d_{\rm long}=2d$ lowers
$V_2$ because Yukawa is convex: $e^{-mx}/x$ is a convex function for
$mx\lesssim1$, so the average over asymmetric bonds gives lower energy
than at the symmetric point.

At $C=0.282$ all pairs are classically bound ($E_0^{(2B)}<0$, see ex29),
so pair interactions drive isosceles deformation.

\subsection{Consequence for quantum Efimov window}

\paragraph{Classical level.}
The true classical ground state is a deformed (isosceles) triangle, not equilateral.
Along the softer mode ($\varepsilon_a=-\varepsilon_b$), the energy decreases
monotonically with deformation at fixed $d$:
\begin{equation}
  E(d,\varepsilon) < E(d,0) \quad \forall\, \varepsilon\neq0.
\end{equation}
The equilateral result $E_{\rm cl}^{\rm eq}=-0.265\,E_{\rm Pl}$ (at $d_{\rm cl}=1.37$)
is an \emph{upper bound} on the true classical minimum.

\paragraph{Quantum level.}
In the quantum problem, the full shape space is the 2D hyperangle manifold.
The equilateral breathing-mode calculation (ex26/ex29) restricts to the
1D submanifold $d_{12}=d_{13}=d_{23}$.  Since the equilateral is a
local energy maximum in the shape directions, the restriction gives an
\emph{upper bound} on $E_0$:
\begin{equation}
  E_0^{\rm equil} \geq E_0^{\rm true}.
\end{equation}
The true ground-state energy is lower (more bound).

\paragraph{Implication for window boundaries.}
A lower true $E_0$ means both thresholds $C_Q^{(3B)}$ and $C_Q^{(2B)}$
are \emph{smaller} than computed in ex26/ex27:
\begin{itemize}
\item $C_Q^{(3B)}$ (lower window bound): shifts downward.
  Binding occurs at weaker coupling.
\item $C_Q^{(2B)}$ (upper window bound): shifts downward.
  Two-body binding also occurs at weaker coupling.
\item The window $m_{\rm sp}\in(0.076,\,0.198)$ computed in ex27 is a
  \emph{conservative estimate}.  The true quantum window is
  \emph{at least as wide} and may extend to slightly larger $m_{\rm sp}$
  (requiring less coupling to bind the triple via the 3-body force).
\end{itemize}

\paragraph{Three-body only vs two-body.}
Whether the 3B-only Efimov window still exists for the full 2D problem
depends on whether $C_Q^{(3B)}$ shifts faster than $C_Q^{(2B)}$ as
the geometry deforms.  Since $V_3$ (the three-body force) is also
symmetric and may have different shape dependence than $V_2$, a full
2D calculation is required.

\subsection{Normal mode frequencies}

The radial (breathing) mode at $d_{\rm cl}=1.37\,l_{\rm Pl}$ has curvature:
\begin{equation}
  \frac{d^2E}{dd^2}\bigg|_{\rm eq} = 0.305\;\text{E}_{\rm Pl}/l_{\rm Pl}^2.
\end{equation}
With effective mass $\mu_r = m_{\rm Pl}/3$, the breathing frequency is
$\omega_r = \sqrt{0.305\times3}\approx0.957\,E_{\rm Pl}$.

The stiffer shape mode has $\lambda_2^{(2)}/d_{\rm cl}^2=+0.047$,
giving $\omega_{\rm sh,2}=\sqrt{0.047\times3}\approx0.375\,E_{\rm Pl}$
(about $0.39\times\omega_r$).
The softer shape mode has $\lambda_1^{(2)}/d_{\rm cl}^2<0$ (imaginary
frequency): this is the unstable mode.

\subsection{Shape stability vs \texorpdfstring{$C$}{C} at quantum mean position}

\begin{table}[h]
\centering
\caption{Shape stability at $d=\langle d\rangle=4.52\,l_{\rm Pl}$
vs coupling $C$.  At $C<0.19$ (near quantum Efimov threshold) the larger
eigenvalue is still positive — saddle-type stability.}
\label{tab:shape_vs_C}
\begin{tabular}{rllll}
\hline
$C$ & $\lambda_1$ & $\lambda_2$ & Type & Note \\
\hline
$0.100$ & $-0.0045$ & $+0.0177$ & Saddle & below 3B threshold \\
$0.128$ & $-0.0074$ & $+0.0123$ & Saddle & classical 3B threshold \\
$0.150$ & $-0.0102$ & $+0.0066$ & Saddle & Efimov window midpoint \\
$0.191$ & $-0.0168$ & $-0.0081$ & Fully unstable & quantum 3B threshold \\
$0.282$ & $-0.0637$ & $-0.0374$ & Fully unstable & $C_{\rm Pl}$ \\
\hline
\end{tabular}
\end{table}

\noindent
The shape mode transitions from ``saddle'' to ``fully unstable'' near
$C\approx0.19$, coinciding with the quantum Efimov threshold
$C_Q^{(3B)}=0.191$.  Below threshold, the stiff shape mode has a
positive restoring force, partially stabilising the equilateral geometry.
Above threshold, both modes are unstable.

\subsection{Summary}

\begin{enumerate}
\item The equilateral triangle is a \textbf{saddle point} in shape space
  for all $C$ and $d$ tested — never a stable minimum.
\item The 1D equilateral breathing-mode calculation (ex26/ex29) gives
  \textbf{upper bounds} on $E_0$ and on the window thresholds.
\item The true quantum ground state is likely a deformed (isosceles)
  triangle.  By $S_3$ symmetry, the quantum state is a superposition
  of three equivalent isosceles configurations, restoring permutation
  symmetry.
\item The quantum Efimov window for $C_{\rm Pl}$ may be slightly wider
  than $(0.076,\,0.198)\,l_{\rm Pl}^{-1}$ (more conservative: could start
  at lower $m_{\rm sp}$ or extend to higher $m_{\rm sp}$).
\item A full 2D hyperangle Schr\"odinger solver is required for
  quantitatively accurate window boundaries.  This is identified as the
  primary open computational problem (ex31+).
\end{enumerate}
